\documentclass[../../../uem_phy_std_q.tex]{subfiles}

\begin{document}
    図のように，真空中において，同形の3枚の金属極板X，Y，Zを%
    平行に狭い間隔で向き合わせてあり，電池$\mathrm{B}_1$（起電力$2V_0$），$\mathrm{B}_2$（起電力$V_0$），%
    スイッチ$\mathrm{S}_1$，$\mathrm{S}_2$およびアースと接続されて回路を形成している．%
    XとYの間隔は$d$，YとZの間隔は$2d$である．%
    XとYからなる平行平板コンデンサの静電容量を$C_0$とし，端の効果は無視する．%
    はじめ，$\mathrm{S}_1$，$\mathrm{S}_2$は共に開いており，いずれの極板も帯電していない．

    \begin{enumerate}
        \item $\mathrm{S}_1$を閉じて安定するまで待った．%
            このとき，X下面の電荷$x_1$，Y下面の電荷$y_1$を求めよ．%
            また，XY間の電場$E_1$（X→Yの向き正），Xの電位$\phi_1$，Yの電位$\theta_1$を求めよ．
        \item $\mathrm{S}_1$を開き，$\mathrm{S}_2$を閉じて安定するまで待った．%
            このとき，X下面の電荷$x_2$，Y下面の電荷$y_2$を求めよ．%
            また，XY間の電場$E_2$（X→Yの向き正），Xの電位$\phi_2$，Yの電位$\theta_2$を求めよ．
        \item $\mathrm{S}_2$を開き，$\mathrm{S}_1$を閉じて安定するまで待った．%
            このとき，X下面の電荷$x_3$，Y下面の電荷$y_3$を求めよ．%
            また，XY間の電場$E_3$（X→Yの向き正），Xの電位$\phi_3$，Yの電位$\theta_3$を求めよ．
        \item ゆっくりとYを移動させ，XとYの間隔が$2d$，YとZの間隔が$d$になるようにした．%
            この間に，$\mathrm{B}_1$を負極→正極の向きに通過した電荷$\Delta Q$を求めよ．
    \end{enumerate}
    \begin{minipage}{\linewidth}
        \centering
            \includegraphics{\subfix{fig_uem_phy_std_em_cb_07/fig_uem_phy_std_em_cb_07_00_q.pdf}}
    \end{minipage}
\end{document}




